% !TeX program = xelatex
% !TeX encoding = UTF-8
\documentclass[12pt,a4paper]{article}
\usepackage[margin=2.5cm]{geometry}
\usepackage{amsmath,amssymb,amsthm,xeCJK,graphicx}
\usepackage[hidelinks]{hyperref}
\usepackage{fancyhdr}
\setCJKmainfont{Microsoft JhengHei}
\makeatletter
\newcommand{\totalpages}{\@abspage@last}
\makeatother
\pagestyle{fancy}
\fancyhf{}
\renewcommand{\headrulewidth}{0pt}
\fancyfoot[C]{\thepage/\totalpages}
\fancypagestyle{plain}{\fancyhf{}\renewcommand{\headrulewidth}{0pt}\fancyfoot[C]{\thepage/\totalpages}}
\setlength{\parindent}{0pt}
\title{Linear System --- Midterm Practice}
\author{葉彥辰\\P46154345}
\date{}
\begin{document}
\maketitle
\begin{center}
\textbf{Mock Examination}\\[6pt]
24 exercises \quad 100 points \quad Suggested time: 240 minutes
\end{center}
\bigskip
\textbf{Instructions}
\begin{itemize}
\item Answer all questions and show your reasoning. Exercises 1--20 carry 4 points each; Exercises 21--24 carry 5 points each.
\item Work by hand. No programming or simulation is required. Exact symbolic answers are preferred.
\item Unless otherwise stated, state variables and inputs are real, $j^2=-1$, and transfer functions use zero initial conditions.
\item Blank solution areas are provided; use extra paper if needed.
\end{itemize}
\bigskip
\textbf{Practice scope}\\
State-space modeling; equilibria and linearization; transfer functions and realizations; linear algebra; eigenvalues and modes; the Cayley--Hamilton theorem; matrix exponentials and impulse responses.
\bigskip

本卷依專案中 HW1--HW5 的原始題目所涵蓋主題重新設計，並非原題逐字彙編，亦非教師公告的正式期中考範圍。僅含題目與空白作答區。
\bigskip

\begin{center}
\begin{tabular}{c|l|r}
Exercises & Topic & Points\\ \hline
1--3 & State-space modeling & 12\\
4--6 & Equilibria and linearization & 12\\
7--11 & Transfer functions and realizations & 20\\
12--17 & Linear algebra and eigenstructure & 24\\
18--20 & Continuous/discrete-time modes & 12\\
21--24 & Matrix functions and responses & 20\\ \hline
 & Total & 100
\end{tabular}
\end{center}
\clearpage
\section*{Exercise 1 \hfill {\normalsize (4 points)}}
Obtain a state-space representation $(A,B,C,D)$ of
\[
q^{(3)}+4\ddot q+5\dot q+2q=u,\qquad y=3q-\dot q+2u.
\]
Define your state variables explicitly.
\subsection*{Solution}
% 在這裡撰寫解答。
\vspace{3cm}
\vfill
\section*{Exercise 2 \hfill {\normalsize (4 points)}}
Obtain a three-state nonlinear state-space model, including the output equation, for
\[
\ddot q_1+\ddot q_2+\sin q_1=u,\qquad
\dot q_2+2q_1+q_2=0,\qquad y=q_1-q_2.
\]
Your state equations must not contain derivatives of the input.
\subsection*{Solution}
% 在這裡撰寫解答。
\vspace{3cm}
\clearpage
\section*{Exercise 3 \hfill {\normalsize (4 points)}}
For $k=0,1,\ldots$, consider
\[
q(k+2)-\tfrac12q(k+1)-\tfrac12q(k)=u(k),\qquad
y(k)=2q(k)-q(k+1)+u(k).
\]
Obtain $(A,B,C,D)$ using $x(k)=[q(k)\;\;q(k+1)]^T$.
\subsection*{Solution}
% 在這裡撰寫解答。
\vspace{3cm}
\vfill
\section*{Exercise 4 \hfill {\normalsize (4 points)}}
Find all equilibrium states and the linearized system matrix at each equilibrium for
\[
\ddot q+2\dot q-q+q^3=0.
\]
Use $x=[q\;\;\dot q]^T$ and define the perturbation state.
\subsection*{Solution}
% 在這裡撰寫解答。
\vspace{3cm}
\clearpage
\section*{Exercise 5 \hfill {\normalsize (4 points)}}
For each scalar system below, find every equilibrium and determine whether a Jacobian linearization exists there. If it exists, compute it; otherwise, justify why it does not.
\[
\text{(a) }\dot x=|x|,\qquad
\text{(b) }\dot x=x|x|,\qquad
\text{(c) }\dot x=\sqrt{|x|}.
\]
\subsection*{Solution}
% 在這裡撰寫解答。
\vspace{3cm}
\vfill
\section*{Exercise 6 \hfill {\normalsize (4 points)}}
Consider the pendulum with torque input
\[
\ddot\theta+0.4\dot\theta+2\sin\theta=u,\qquad y=\sin\theta.
\]
For $u^e=1$, find all equilibrium states. Obtain $(A,B,C,D)$ for the perturbation model about each equilibrium, using $x=[\theta\;\;\dot\theta]^T$.
\subsection*{Solution}
% 在這裡撰寫解答。
\vspace{3cm}
\clearpage
\section*{Exercise 7 \hfill {\normalsize (4 points)}}
Under zero initial conditions, find the transfer matrix from $[u_1\;\;u_2]^T$ to $[y_1\;\;y_2]^T$:
\[
\dot y_1+2y_1-y_2=u_1,\qquad
\dot y_2+3y_2=\dot u_2+u_1.
\]
State the dimensions of the transfer matrix.
\subsection*{Solution}
% 在這裡撰寫解答。
\vspace{3cm}
\vfill
\section*{Exercise 8 \hfill {\normalsize (4 points)}}
Let $h>0$ and
\[
\dot x(t)=-3x(t)+x(t-h)+u(t),\qquad y(t)=2x(t),
\]
with zero prehistory $x(t)=0$ for $t\leq0$. Find $G(s)=Y(s)/U(s)$. Explain whether $G(s)$ can be realized exactly by a finite-dimensional continuous-time LTI model with constant matrices and no delays.
\subsection*{Solution}
% 在這裡撰寫解答。
\vspace{3cm}
\clearpage
\section*{Exercise 9 \hfill {\normalsize (4 points)}}
Obtain a controllable companion-form realization of
\[
G(s)=\frac{2s^2+7s+5}{s^2+4s+3}.
\]
Use the displayed denominator without cancelling factors first. Then determine whether your realization is minimal and, if needed, give a lower-order realization.
\subsection*{Solution}
% 在這裡撰寫解答。
\vspace{3cm}
\vfill
\section*{Exercise 10 \hfill {\normalsize (4 points)}}
Construct a real state-space realization of
\[
G(s)=\begin{bmatrix}
\dfrac1{s+2}&\dfrac1{s+2}\\[4pt]
\dfrac1{s+4}&-\dfrac1{s+4}
\end{bmatrix}.
\]
Use no more than two states. Verify your answer by computing $C(sI-A)^{-1}B+D$.
\subsection*{Solution}
% 在這裡撰寫解答。
\vspace{3cm}
\clearpage
\section*{Exercise 11 \hfill {\normalsize (4 points)}}
Consider
\[
A=\begin{bmatrix}0&1\\-2&-3\end{bmatrix},\quad
B=\begin{bmatrix}0\\1\end{bmatrix},\quad C=\begin{bmatrix}1&0\end{bmatrix},\quad D=0.
\]
With $z=Tx$ and $T=\begin{bmatrix}1&1\\0&1\end{bmatrix}$, compute the transformed matrices. Show that the original and transformed realizations have the same transfer function.
\subsection*{Solution}
% 在這裡撰寫解答。
\vspace{3cm}
\vfill
\section*{Exercise 12 \hfill {\normalsize (4 points)}}
Determine all real values of $a,b$ for which the following equations are consistent:
\[
x_1+x_2+x_3=1,\qquad
2x_1+2x_2+2x_3=a,\qquad
x_1+x_2+bx_3=2.
\]
For each consistent case, describe all solutions and state whether the solution is unique.
\subsection*{Solution}
% 在這裡撰寫解答。
\vspace{3cm}
\clearpage
\section*{Exercise 13 \hfill {\normalsize (4 points)}}
For
\[
A=\begin{bmatrix}1&2&0&1\\0&1&1&1\\1&3&1&2\end{bmatrix},
\]
find bases for $\ker A$ and $\operatorname{range}A$. Compute the rank and nullity and verify the rank--nullity identity.
\subsection*{Solution}
% 在這裡撰寫解答。
\vspace{3cm}
\vfill
\section*{Exercise 14 \hfill {\normalsize (4 points)}}
Let
\[
v_1=\begin{bmatrix}1\\0\\1\end{bmatrix},\quad
v_2=\begin{bmatrix}0\\1\\1\end{bmatrix},\quad
v_3=\begin{bmatrix}1\\1\\2\end{bmatrix}.
\]
Determine the condition on $(a,b,c)$ for $[a\;\;b\;\;c]^T$ to belong to their span. When that condition holds, give all coefficient triples in the linear combination. Are $v_1,v_2,v_3$ linearly independent?
\subsection*{Solution}
% 在這裡撰寫解答。
\vspace{3cm}
\clearpage
\section*{Exercise 15 \hfill {\normalsize (4 points)}}
Find the eigenvalues and a basis for each eigenspace of
\[
A=\begin{bmatrix}2&1&0\\0&2&0\\0&0&-1\end{bmatrix}.
\]
State the algebraic and geometric multiplicities of each eigenvalue, and determine whether $A$ is diagonalizable over $\mathbb R$.
\subsection*{Solution}
% 在這裡撰寫解答。
\vspace{3cm}
\vfill
\section*{Exercise 16 \hfill {\normalsize (4 points)}}
Construct a real $3\times3$ companion matrix, with ones on the superdiagonal, whose eigenvalues are
\[
-2,\qquad -1+2j,\qquad -1-2j,\qquad j^2=-1.
\]
Verify its characteristic polynomial.
\subsection*{Solution}
% 在這裡撰寫解答。
\vspace{3cm}
\clearpage
\section*{Exercise 17 \hfill {\normalsize (4 points)}}
Let $v,w\in\mathbb R^4$ be nonzero and let $A=vw^T$. Find the rank and nullity of $A$. Determine all eigenvalues, with algebraic multiplicities, and whether $A$ is diagonalizable in each case:
\[
\text{(a) }w^Tv=3,\qquad \text{(b) }w^Tv=0.
\]
Justify your conclusions.
\subsection*{Solution}
% 在這裡撰寫解答。
\vspace{3cm}
\vfill
\section*{Exercise 18 \hfill {\normalsize (4 points)}}
A real $3\times3$ matrix $A$ has eigenvectors
\[
v_1=\begin{bmatrix}1\\0\\0\end{bmatrix},\quad
v_2=\begin{bmatrix}0\\1\\1\end{bmatrix},\quad
v_3=\begin{bmatrix}0\\1\\-1\end{bmatrix}
\]
corresponding to $-2,-1,1$, respectively. Find the response of $\dot x=Ax$ with $x(0)=[2\;\;3\;\;1]^T$. Characterize all initial states whose responses tend to zero as $t\to\infty$.
\subsection*{Solution}
% 在這裡撰寫解答。
\vspace{3cm}
\clearpage
\section*{Exercise 19 \hfill {\normalsize (4 points)}}
Let $A\in\mathbb R^{2\times2}$ satisfy
\[
A\begin{bmatrix}1\\j\end{bmatrix}
=(-1+2j)\begin{bmatrix}1\\j\end{bmatrix},\qquad j^2=-1.
\]
Find the real response of $\dot x=Ax$ for $x(0)=[2\;\;-1]^T$. Determine whether this nonzero response is periodic and whether it tends to zero.
\subsection*{Solution}
% 在這裡撰寫解答。
\vspace{3cm}
\vfill
\section*{Exercise 20 \hfill {\normalsize (4 points)}}
For the discrete-time system
\[
x(k+1)=Ax(k),\qquad
A=\begin{bmatrix}0&-\tfrac12\\[2pt]\tfrac12&0\end{bmatrix},\qquad
x(0)=\begin{bmatrix}1\\0\end{bmatrix},
\]
find a real closed-form expression for $x(k)$. Is the response periodic? Does it tend to zero? Explain both answers.
\subsection*{Solution}
% 在這裡撰寫解答。
\vspace{3cm}
\clearpage
\section*{Exercise 21 \hfill {\normalsize (5 points)}}
Suppose $A\in\mathbb R^{3\times3}$ and
\[
\det(sI-A)=s^3+4s^2+5s+2.
\]
Using the Cayley--Hamilton theorem, express (a) $A^3$, (b) $A^5$, and (c) $A^{-1}$ as linear combinations of $I,A,A^2$. Justify the existence of $A^{-1}$.
\subsection*{Solution}
% 在這裡撰寫解答。
\vspace{3cm}
\vfill
\section*{Exercise 22 \hfill {\normalsize (5 points)}}
For
\[
A=\begin{bmatrix}-2&1\\0&-2\end{bmatrix},
\]
compute $e^{At}$ in two ways: (a) using the nilpotent part of $A$ and the exponential series; (b) by inverse Laplace transformation of $(sI-A)^{-1}$. Use your result to find $x(t)$ for $\dot x=Ax$, $x(0)=[0\;\;1]^T$.
\subsection*{Solution}
% 在這裡撰寫解答。
\vspace{3cm}
\clearpage
\section*{Exercise 23 \hfill {\normalsize (5 points)}}
Consider the initially resting system
\[
\dot x=\begin{bmatrix}-1&0\\0&-3\end{bmatrix}x+
\begin{bmatrix}1\\1\end{bmatrix}u,\qquad
y=\begin{bmatrix}1&0\end{bmatrix}x,\qquad x(0^-)=0.
\]
For $u(t)=\delta(t)$, find $x(0^+)$, $x(t)$ for $t>0$, and $y(t)$. Determine which system modes occur in the state response and in the output response, and explain the difference.
\subsection*{Solution}
% 在這裡撰寫解答。
\vspace{3cm}
\vfill
\section*{Exercise 24 \hfill {\normalsize (5 points)}}
Let $A\in\mathbb R^{n\times n}$ have eigenvalue $j\omega$, where $\omega>0$, with eigenvector $v=p+jq$, $p,q\in\mathbb R^n$.
Prove that $\dot x=Ax$ has a nonzero real periodic solution. Construct a suitable real initial state from $p,q$ and give a period. Does the existence of this solution imply that every solution is bounded? Justify your answer with a proof or counterexample.
\subsection*{Solution}
% 在這裡撰寫解答。
\vspace{3cm}
\clearpage
\end{document}
